% ============================================================================
%  Section 1 — Introduction
%
%  Mark API examples are shown throughout this section to demonstrate
%  the marking workflow. Remove/replace them with your actual content.
% ============================================================================
\section{\Mark{changes}[heading]{Introduction}}

\Mark{changes}{%
This is an example of a highlighted \emph{change} — a new or revised sentence
is wrapped in \texttt{\textbackslash Mark\{changes\}\{...\}} to highlight it
in the compiled PDF. The highlight color (yellow by default) is defined in
\texttt{config.tex} via \texttt{\textbackslash DefineMarkRole\{changes\}\{\textbackslash ChangesColor\}}.%
}

\Mark{missing}{TODO: Write the introduction here. %
Motivate the problem, state the main contribution, and outline the paper structure.}

Here is an example of a \emph{missing} equation placeholder:

\MarkPlaceholderFigure{missing}{%
  Add a system diagram here showing the high-level architecture.%
}{fig:system}{System overview.}

Here is an example of a highlighted block equation:

\begin{MarkEnv}{changes}[math]
    \mathcal{F}(\alpha) =
    \sum_{k=1}^{K}
    \alpha \cdot f_k + (1 - \alpha) \cdot g_k
\end{MarkEnv}

\noindent
where \Mark{changes}{$\alpha \in [0, 1]$ is a risk parameter} %
and \Mark{missing}{TODO: define $f_k$ and $g_k$}.

The paper is organized as follows.
\Cref{sec:background} \Mark{missing}{TODO: describe each section briefly.}

% ============================================================================
%  Related Work (can be a subsection of Introduction or a standalone section)
% ============================================================================
\subsection{Related Work}

\Mark{missing}{TODO: Survey related work here. %
Cite key papers from the literature knowledge base using \texttt{@find-papers}.}

Prior work on \Mark{changes}{relevant topic}~\cite{example2024} has shown that
\Mark{missing}{TODO: summarize the key finding and its limitation.}
